% 30_body.tex — 산출물 ㉚ 시드 텀시트 · 캡테이블 v1 (빈 템플릿 / 완성 예시 공용 본문) — gen_product_templates.py 가 Product_specs.py 에서 생성
% 드라이버: 30_시드텀시트캡테이블_빈템플릿.tex (\wbblanktrue) / 30_시드텀시트캡테이블_완성예시.tex (\wbblankfalse — 워크북 §X.5 확정 후)
\documentclass[10pt,a4paper]{article}
\usepackage[a4paper,margin=16mm,top=14mm,bottom=20mm]{geometry}
\def\DocNum{30}\def\DocDay{6}\def\DocIdx{5}
\def\DocTitle{시드 텀시트 · 캡테이블 v1}
\def\DocSub{Seed Term Sheet \& Cap Table v1}
\def\DocVariant{\B{빈 템플릿}{딥게이지 완성 예시}}
\usepackage{wbstyle}
\def\DocCirc{㉚}   % newunicodechar(㉛~㊵ 폴백) 뒤에 정의해야 활성 문자로 읽힌다
\begin{document}
\docheader
 {\Bg{[회사명]}{주식회사 딥게이지}}
 {\Bg{[v1.x / D+n]}{v1.0 / 2026-08-21 (D+172) · 부속 메모 D+186}}
 {\Bg{[작성자 — CEO]}{강민수\rd{7mm}}}
 {\Bg{[검토자 — 법무·기존 주주]}{법무(외부) · 오세진 · 윤하경\rd{7mm}}}

%% ------------------------------------------------------------------ 1
\secbar{1. 조달 조건 요약}
\guide{RCPS 12억 / pre 48 · post 60 / 지분 20\% / 참가적 1x / 상환권 연 3\% / drag 60\% / TIPS 8억(조건: 6개월 내 유료 5개사). 각 조건이 우리에게 무슨 뜻인지 한 줄.}
{\footnotesize
\begin{tabularx}{\linewidth}{|L{40mm}|L{34mm}|Y|}\hline
\hd{조건} & \hd{값} & \hd{우리에게 무슨 뜻인가} \\ \hline
\ifwbblank
\rd{9mm}투자 금액 / 증권 종류 &  &  \\ \hline
\rd{9mm}pre / post-money &  &  \\ \hline
\rd{9mm}지분(FD) &  &  \\ \hline
\rd{9mm}청산우선권(배수·참가 여부) &  &  \\ \hline
\rd{9mm}상환권 &  &  \\ \hline
\rd{9mm}drag-along 발동 요건 &  &  \\ \hline
\rd{9mm}이사회·정보권 &  &  \\ \hline
\rd{9mm}TIPS 연계 조건 &  &  \\ \hline
\else
투자 금액 / 증권 종류 & 12억 원 / RCPS(선투자 2억 D+150 포함) & 12명 기준 17개월. D+149 잔액 1,505만 — 없으면 D+156에 급여가 없다 \\ \hline
pre / post-money & pre 48억 · post 60억 & 주당 4,800원(48억 ÷ 1,000,000주). '회사 가치'가 아니라 이 라운드의 가격 \\ \hline
지분(FD) & 노들 20.0\% & 250,000주 신주. 창업 3인 72\%(45 → 36, 30 → 24, 15 → 12), 풀 8\% \\ \hline
청산우선권(배수·참가 여부) & 1x 참가적, 상한 없음 & 12억을 먼저 받고 나머지의 20\%를 또 받는다. 100억 매각에서 29.6억(비참가적이면 20억) — 항목 3 \\ \hline
상환권 & 3년 후, 연 3\% 복리, 배당 가능 이익 한도 & 3년 뒤 13.1억을 요구할 수 있으나 이익잉여금 한도 — 실질은 협상 압력 \\ \hline
drag-along 발동 요건 & FD 60\%, 투자자 동의 조항 없음 & 노들 20 + 강민수 36 = 56\% < 60\% — 창업자 둘 이상 필요. 창업 3인 72\%만으로 발동 가능(3번 협상) \\ \hline
이사회·정보권 & 3인(강민수·오세진 + 노들 1) + 옵저버 · 월간 보고 & 과반은 창업자. 월간 보고의 ARR 정의는 우리가 정한다(㉘ 항목 6) \\ \hline
TIPS 연계 조건 & 운영사 추천 → 일반트랙 8억(2년) · 추천 조건 D+360까지 유료 5개사 & 8억은 R\&D 자금(희석 없음). 미달 시 후속 추천·2차년도 — 항목 4 \\ \hline
옵션풀 위치(pre/post) & 조항 없음 — D+186 부속 메모 & ㉑ 항목 3 요구가 협상 1번, 구두 합의, 문서에 없다. 없는 조항은 다음 리드가 정한다 \\ \hline
\fi
\end{tabularx}}

%% ------------------------------------------------------------------ 2
\secbar{2. 캡테이블 T0 → T1}
\guide{주식수·지분(FD)을 라운드 전후로. 옵션풀이 어디서 나오는지(pre면 기존 주주 희석) 명시.}
{\footnotesize
\begin{tabularx}{\linewidth}{|L{34mm}|L{24mm}|L{18mm}|L{24mm}|L{18mm}|Y|}\hline
\hd{주주} & \hd{T0 주식수} & \hd{T0 지분} & \hd{T1 주식수} & \hd{T1 지분} & \hd{비고} \\ \hline
\ifwbblank
\rd{8mm}강민수 &  &  &  &  &  \\ \hline
\rd{8mm}오세진 &  &  &  &  &  \\ \hline
\rd{8mm}윤하경 &  &  &  &  &  \\ \hline
\rd{8mm}옵션풀 &  &  &  &  &  \\ \hline
\rd{8mm}노들투자파트너스(RCPS) &  &  &  &  &  \\ \hline
\rd{8mm}계(FD) &  &  &  &  &  \\ \hline
\else
강민수 & 450,000 & 45.0\% & 450,000 & 36.0\% & 주식수 불변 — 희석 \\ \hline
오세진 & 300,000 & 30.0\% & 300,000 & 24.0\% &  \\ \hline
윤하경 & 150,000 & 15.0\% & 150,000 & 12.0\% &  \\ \hline
옵션풀 & 100,000 & 10.0\% & 100,000 & 8.0\% & 추가 풀 없음. 위치 미정 — 시리즈A 15\% 요구 시 정해진다 \\ \hline
노들투자파트너스(RCPS) & — & — & 250,000 & 20.0\% & 12억 ÷ 4,800원 \\ \hline
계(FD) & 1,000,000 & 100\% & 1,250,000 & 100\% &  \\ \hline
\fi
\end{tabularx}}

%% ------------------------------------------------------------------ 3
\secbar{3. 매각가 시나리오 — 참가적/비참가적 × 옵션풀 pre/post × 매각가 3구간}
\guide{매각가 30억 / 100억 / 300억에서 투자자와 창업자가 각각 얼마를 받는가. 참가적 1x가 어느 구간에서 얼마나 가져가는지 손으로 계산한다.}
{\footnotesize
\begin{tabularx}{\linewidth}{|L{44mm}|Y|Y|Y|}\hline
\hd{시나리오} & \hd{30억} & \hd{100억} & \hd{300억} \\ \hline
\ifwbblank
\rd{9mm}비참가적 · 풀 post &  &  &  \\ \hline
\rd{9mm}비참가적 · 풀 pre &  &  &  \\ \hline
\rd{9mm}참가적 1x · 풀 post &  &  &  \\ \hline
\rd{9mm}참가적 1x · 풀 pre &  &  &  \\ \hline
\rd{9mm}창업자 3인 합계(참가적·pre) &  &  &  \\ \hline
\else
비참가적 · 풀 post(투자자 19.6\%) & 12.0 / 15.8 & 19.6 / 70.4 & 58.7 / 211.3 \\ \hline
비참가적 · 풀 pre(20.0\%) & 12.0 / 15.7 & 20.0 / 70.0 & 60.0 / 210.0 \\ \hline
참가적 1x · 풀 post & 15.5 / 12.7 & 29.2 / 62.0 & 68.3 / 202.9 \\ \hline
참가적 1x · 풀 pre & 15.6 / 12.6 & 29.6 / 61.6 & 69.6 / 201.6 \\ \hline
창업자 3인 합계(참가적·pre) & 12.6 & 61.6 & 201.6 \\ \hline
(참고) 실제 K-02 — 추가 풀 없음, 참가적 & 15.6 / 13.0 & 29.6 / 63.4 & 69.6 / 207.4 \\ \hline
\fi
\end{tabularx}}

%% ------------------------------------------------------------------ 4
\secbar{4. TIPS 조건과 제품 계획의 충돌}
\guide{'6개월 내 유료 5개사'가 ㉖ 범위 절단·㉘ 가격에 미치는 압력. 무엇을 앞당기고 무엇을 포기하는가.}
{\footnotesize
\begin{tabularx}{\linewidth}{|L{40mm}|Y|Y|}\hline
\hd{조건} & \hd{제품 계획에의 압력} & \hd{대응} \\ \hline
\ifwbblank
\rd{11mm} &  &  \\ \hline
\rd{11mm} &  &  \\ \hline
\rd{11mm} &  &  \\ \hline
\else
D+360까지 유료 5개사 & 라인 1개·6개월 계약을 받고 싶어진다 → ACV 반토막, 실사용 데이터 부족 & 최소 라인 2·12개월 유지(㉘ 항목 7). 파이프라인 4사 알파 라인으로 · 세일즈 채용 D+200 \\ \hline
같은 조건 & 온프렘 요구 고객(H사 계열)을 넣고 싶어진다 → Won't 붕괴, 6.0 mm & ㉖ 항목 5 유지. 도면 미업로드 모드로 풀리는 고객 먼저(B-19) \\ \hline
2차년도 3억은 1차년도 평가 후 & 평가 지표를 '유료 고객 수'로만 보고하고 싶어진다 & 월간 보고에 계약/실사용 라인(G-18)과 오류 예산 소진 여부 병기 — 서영채에게 먼저 \\ \hline
\fi
\end{tabularx}}

%% ------------------------------------------------------------------ 5
\secbar{5. 협상 요청 3개와 BATNA}
\guide{바꾸고 싶은 조항 3개(우선순위·근거·대안 문구)와 결렬 시 대안(다른 투자자·매출로 버티기·규모 축소).}
{\footnotesize
\begin{tabularx}{\linewidth}{|L{12mm}|L{34mm}|Y|Y|}\hline
\hd{순위} & \hd{조항} & \hd{요청 문구} & \hd{근거} \\ \hline
\ifwbblank
\rd{11mm} &  &  &  \\ \hline
\rd{11mm} &  &  &  \\ \hline
\rd{11mm} &  &  &  \\ \hline
\else
1 & 옵션풀 위치 & '본 라운드 이후 옵션풀 신설·증액 시 그 희석은 라운드 후 전체 주주가 지분 비율대로 부담한다.' & ㉑ 항목 3. pre면 0.43\%p 이전(항목 3 가상 비교). 결과: 구두 합의 → 미기재 — D+186 발견 \\ \hline
2 & 청산우선권 & '참가적 배당은 원금의 3배를 상한으로 한다' 또는 '비참가적 1x' & 100억에서 9.6억 차이. 결과: 거절 — '시드 RCPS 표준'. 상한 없이 서명 \\ \hline
3 & drag-along & '발동 요건 60\%에 창업자 과반의 동의를 포함한다' & 후속 라운드에서 투자자 합계 40\% 초과 시 창업자 1인만 붙으면 발동. 결과: 수용 \\ \hline
\fi
\end{tabularx}}

%% ------------------------------------------------------------------ 6
\secbar{6. 서명·다음 라운드 전제}
\guide{다음 라운드(브릿지·시리즈A)에 이 조건들이 어떻게 이월되는지 한 줄 후 서명.}
{\footnotesize
\begin{tabularx}{\linewidth}{|Y|Y|Y|Y|}\hline
\hd{강민수 CEO} & \hd{오세진 CTO} & \hd{윤하경 CPO} & \hd{(투자자 대표)} \\ \hline
\ifwbblank
\rd{16mm} & & &  \\ \hline
\else
\rd{16mm} 강민수 (서명) & 오세진 (서명) & 윤하경 (서명) & [노들] 서영채 (서명) \\ \hline
\fi
\rd{7mm}일자 & & &  \\ \hline
\end{tabularx}}

\end{document}